% ------------------------------------------------------------------------
% -*-TeX-*- -*-Hard-*- Smart Wrapping
% ------------------------------------------------------------------------


\chapter{関連研究}
\label{cha:2}

本章では，本研究に関連する研究を4つの系統——（1）マルチエージェント議論，（2）モデルマージとLoRA合成の理論，（3）進化的アプローチによるモデル最適化，（4）協調の定量化と適応度設計の理論——に整理し，最後に本研究の位置づけを述べる．

\section{マルチエージェント議論}
\label{sec:2.1}

\subsection{原典と初期の成果}
\label{sec:2.1.1}

マルチエージェント議論（Multi-Agent Debate; MAD）の原典とされるのがDuら\cite{du2023improving}の研究である．複数のLLMインスタンスがまず独立に回答を生成し，続くラウンドで他エージェントの回答を参照しながら自身の回答を批判的に更新する，という単純な手続きにより，3エージェント$\times$2ラウンドの構成でGSM8Kを77から85へ，MMLUを63.9から71.1へ改善できることが示された．同研究では，他者回答を要約して提示することが有効である一方，単一エージェントの自己反省のみでは改善が得られないことも報告されており，「他者からの異なる情報」が改善の源泉であることが示唆されている．並行してLiangら\cite{liang2023encouraging}は，単一モデルの自己反省が自身の初期解に固執する「思考の退化（Degeneration-of-Thought; DoT）」を指摘し，異なる立場を割り当てられたエージェント間の対立的議論がDoTの回避に必要であると論じた．この知見は，議論に参加するエージェントに明示的な視点の多様性——本研究の文脈ではペルソナの多様性——を与えるべきであるという本研究の設計動機の直接の源流である．

議論を実行時のプロトコルとしてだけでなく訓練信号の源として使う方向も発展した．Subramaniamら\cite{subramaniam2025multiagent}は，議論の軌跡から生成役と批評役を分化させて自己教師あり微調整（SFT）を行うMultiagent Finetuningを提案し，Phi-3のMATH精度を5反復で58.8から66.0へ改善するとともに，エージェント間の多様性の維持が反復改善の持続と相関することを示した．またMACA\cite{maca2025preference}は議論トレースからの選好学習により小型モデルのGSM8Kを$+27.6$\%改善した．さらにParkら\cite{park2025maporl}のMAPoRLは，議論における協調品質を報酬としてマルチエージェント強化学習で複数モデルを共訓練する枠組みを提案しており，「協調の質を最適化信号にする」という点で本研究に最も思想が近い．ただしMAPoRLは勾配ベースのRLであり，重み空間の集団に対する世代交代型の進化ではない（この対比は\ref{sec:2.5}節で詳述する）．

このほか，複数モデルの出力を階層的に統合するMixture-of-Agents\cite{wang2024mixture}に対して，Liら\cite{li2025rethinking}は単一の強いモデルの複数サンプルを統合するSelf-MoAが異種混合を上回る場合が多いことを示し，混合が有効なのは構成モデルの品質が拮抗している場合に限られると分析した．品質と多様性はトレードオフの関係にあり，多様性は品質を犠牲にしない範囲でのみ価値を持つ——この知見は，品質（貢献度）と多様性を単一の適応度に統合する本研究の設計課題を先取りしている．

\subsection{批判研究：MADは本当に有効か}
\label{sec:2.1.2}

MADの初期の成功報告に対しては，その後，系統的な批判が積み重ねられている．第一に，Smitら\cite{smit2024should}は複数のMAD戦略を計算量を統制して比較し，MADがSelf-Consistency（複数サンプルの多数決）\cite{wang2023selfconsistency}等の単純な手法に安定して勝てないことを示した．これはMADの効果検証において計算量を揃えたベースラインが不可欠であることを意味する．第二に，Zhangら\cite{zhang2025debate}は既存MAD研究の評価設計（ベースラインの弱さ，ハイパーパラメータの非統制等）を系統的に指摘し，公平な条件ではMADの優位が大きく縮小すること，そして今後の方向としてエージェント間の異質性の活用を提案した．第三に，議論という相互作用そのものの失敗様式も特定されている．他者への追従（sycophancy）により，当初正答を持っていたエージェントが誤答へ引きずられる「正$\rightarrow$誤の伝染」が観察されており\cite{talk2025cheap}，7〜8B級の同一モデルの複製からなる同質なMADは非効率かつ不安定であることが報告されている\cite{cost2026consensus}．

\subsection{成功条件の知見と本研究への含意}
\label{sec:2.1.3}

批判研究と並行して，「どのような条件なら議論が機能するのか」を切り分ける研究も現れている．Choiら\cite{choi2025debate}はMADの利得の大半が最終的な多数決（投票）で説明でき，議論中の信念更新はマルチンゲール的（期待値的に情報を増やさない）であることを示した．Kaesbergら\cite{kaesberg2025voting}は集約方式を比較し，推論タスクでは合意形成よりも投票が優れ（$+13.2$\%），議論ラウンドの追加はむしろ逆効果であることを示した．さらに近年の分析は，初期回答の多様性と確信度に条件付けられた回答更新がMAD成功の鍵であることを示している\cite{demystifying2026mad}．

これらを総合すると，次の緊張関係が浮かび上がる．\textbf{議論の利得は「多様な初期回答の生成」と「正しい方向への選択的更新」から生じるが，同質で追従的な小型モデル集団はそのどちらも満たしにくい．}
7〜8B級ですら同質MADの不安定性が報告されている以上，本研究が対象とする4B級では，素の（無調整の）debateは機能しない可能性が高い．本研究はこの緊張関係を回避すべき前提ではなく検証対象として引き受ける．すなわち，（i）プロトコル面では文献の成功条件に忠実に従い（独立初期回答，最大2ラウンド，多数決集約），（ii）エージェント面では重み空間に焼き込まれたペルソナ多様性とチームレベル選択圧によって「議論が機能するエージェント集団」を明示的に育成できるかを問う．

\section{モデルマージとLoRA合成の理論}
\label{sec:2.2}

本研究の進化アルゴリズムはLoRA重みの交叉（ブレンド）を核とするため，重み合成がいつ・なぜ妥当なのかという理論的根拠を必要とする．

重み平均によるモデル合成の原点はModel Soups\cite{wortsman2022model}である．同一の事前学習モデルから微調整された複数モデルの重みを単純平均するだけで精度と頑健性が向上することが示されたが，これは微調整解が損失地形上で線形に接続されている（linear mode connectivity）ことを前提とする．Ilharcoら\cite{ilharco2022editing}のTask Arithmeticは，微調整前後の重み差分$\Delta W$を「タスクベクトル」とみなし，その加減算によって能力を合成・除去できることを示し，$\Delta W$空間での算術という枠組みを確立した．

一方で，素朴な線形和には劣化要因があることも明らかになっている．Yadavら\cite{yadav2023ties}のTIES-Mergingは，タスクベクトル間のパラメータ符号の衝突（干渉）が線形和の性能劣化の主因であることを特定し，符号衝突を解決してからマージする手法を提案した．Yuら\cite{yu2023language}のDAREは，$\Delta W$の大部分をランダムにdropして残りをrescaleしてもタスク性能が保たれる冗長性を利用し，干渉を減らしてマージする．

LoRAに固有の問題を扱ったのがStoicaら\cite{stoica2024knots}のKnOTSである．同研究は，LoRAによって得られた$\Delta W$はフルファインチューニングの場合よりもモデル間の整合（alignment）が低く，素朴な重み平均が顕著に劣化することを示し，複数の$\Delta W$をSVDによって共通基底に整列させてからマージする手法を提案した．さらに，LoRAの低ランク因子$A$，$B$を別々に補間する方式には\textbf{交差項問題}がある．すなわち，2つのアダプタ$(B_1 A_1)$，$(B_2 A_2)$の因子を係数$\alpha$で別々に補間すると，
\begin{equation}
\bigl((1-\alpha)B_1 + \alpha B_2\bigr)\bigl((1-\alpha)A_1 + \alpha A_2\bigr)
= (1-\alpha)^2 B_1 A_1 + \alpha^2 B_2 A_2 + \alpha(1-\alpha)\bigl(B_1 A_2 + B_2 A_1\bigr)
\end{equation}
となり，意図した凸結合$(1-\alpha)B_1 A_1 + \alpha B_2 A_2$ではなく，意味を持たない交差項$B_1 A_2$，$B_2 A_1$が混入する．加えて$BA = (BR)(R^{-1}A)$という再パラメータ化不変性のため，同一の$\Delta W$を表す因子分解は無数に存在し，因子空間での補間は分解の取り方に依存してしまう．したがって，LoRA集団の交叉演算子は因子空間ではなく$\Delta W$空間で定義し，必要に応じてSVDで低ランクに再分解するのが理論的に正当である．本研究の交叉演算子（$\Delta W$空間ブレンド＋SVD再分解）はこの知見に直接依拠しており，因子別補間（naive方式）はアブレーション対象（A3）として実装内に位置づける（実施状況は\S\ref{sec:4.3}）．

なお，線形マージ以外の合成として，Prabhakarら\cite{prabhakar2024lora}は複数LoRAをランク方向に連結するCAT（ランクの合計をとる連結）が線形マージを上回り，数学タスクで$+43$\%の改善を得ることを報告している．モデルマージ全般の体系的整理はYangら\cite{yang2024model}のサーベイに詳しい．

\section{進化的アプローチによるモデル最適化}
\label{sec:2.3}

進化計算をLLMの重み・合成レシピの探索に適用する研究は，近年ひとつの系統を成しつつある．

先駆けとなったのはAkibaら（Sakana AI）の進化的モデルマージである\cite{akiba2024evolutionary}．CMA-ESを用いて既存フルモデル群のマージレシピ（層別のマージ係数や推論経路）を探索し，日本語LLMであるEvoLLM-JP 7Bが一部ベンチマークで70B級モデルを超える性能を達成した．ただしこれは既存モデルを素材とする「レシピの1点探索」であり，モデル集団自体の世代交代は行われない．LoraHub\cite{huang2023lorahub}も同様に，既存LoRA群の合成係数ベクトルを勾配フリー最適化で求める1点解の探索であり，かつ$A$・$B$因子別の合成であるため前節の交差項問題を内包する．

集団ベースの探索へ進んだのがModel Swarms\cite{feng2024model}である．複数のLoRAエキスパートを粒子とみなし，粒子群最適化（PSO）によって重み空間を協調探索することで，200例程度の少数データから最大$+21$\%の適応改善を得た．ただしPSOは速度ベクトルによる連続的な移動であり，交叉・突然変異・世代交代という遺伝的操作を持たず，効用関数も個体単位である．

遺伝的アルゴリズムをLoRA重みに全面的に適用したのがGENOME/GENOME+である\cite{zhang2025genome}．LoRA重みを遺伝子とみなし，交叉・突然変異・選択・継承（succession）からなる世代交代を最大40個体の集団で実行する．本研究にとって最近接の先行研究であるが，その適応度は\textbf{個体単独のタスク性能のみ}である．同系統のEvoPref\cite{evopref2026evolutionary}は32個体のLoRAをNSGA-IIで200世代進化させ多様性アーカイブを併用するが，目的はアライメントであり評価は個体単位である．またPopuLoRA\cite{creuscastanyer2026populora}は共有ベースモデル上のLoRA集団を出題者対解答者の敵対的self-playで共進化させ，同ランク制約下の交叉演算子を提案した．個体間の相互作用が適応度を決めるという点で本研究に近づいているが，その相互作用は敵対的であり，協調的な議論への寄与を測るものではない．実装基盤としては，進化的マージの汎用ライブラリMergenetic\cite{minut2025mergenetic}が公開されている．

主要研究の対比を表\ref{tab:2.1}に示す．

\begin{table}[htbp]
\caption{進化的モデル最適化の主要研究の対比}
\label{tab:2.1}
\centering
\footnotesize
\setlength{\tabcolsep}{3pt}
\begin{tabular}{|p{2.5cm}|p{2.3cm}|p{1.9cm}|p{2.3cm}|p{2.4cm}|p{1.9cm}|}
\hline
研究 & 進化の対象 & 探索方式 & 集団と世代交代 & 適応度の評価単位 & 個体間相互作用 \\
\hline
進化的モデルマージ\cite{akiba2024evolutionary} & マージレシピ（層別係数・推論経路） & CMA-ES & なし（レシピ1点解） & マージ後モデル単体の性能 & なし \\
\hline
LoraHub\cite{huang2023lorahub} & LoRA合成係数ベクトル & 勾配フリー最適化 & なし（1点解） & 合成モデル単体の性能 & なし \\
\hline
Model Swarms\cite{feng2024model} & LoRA重み（粒子） & PSO & 集団あり・世代交代なし & 個体単位の効用 & 探索情報の共有のみ \\
\hline
GENOME/ GENOME+\cite{zhang2025genome} & LoRA重み（遺伝子） & GA（交叉・突然変異・選択） & 集団あり・世代交代あり（最大40個体） & 個体単独のタスク性能 & なし \\
\hline
EvoPref\cite{evopref2026evolutionary} & LoRA重み & NSGA-II（200世代） & 集団あり・世代交代あり（32個体） & 個体単位（アライメント目的） & なし \\
\hline
PopuLoRA\cite{creuscastanyer2026populora} & LoRA重み & 共進化GA & 集団あり・世代交代あり & 対戦の勝敗 & \textbf{敵対的} self-play \\
\hline
\textbf{本研究} & LoRA重み（役割別サブ集団） & 協調的共進化GA & 集団あり・世代交代あり & \textbf{チーム議論への厳密Shapley寄与} & \textbf{協調的}議論 \\
\hline
\end{tabular}
\end{table}

表\ref{tab:2.1}が示すとおり，系統の発展は「レシピの1点探索 $\rightarrow$ 集団による重み空間探索 $\rightarrow$ 世代交代型の遺伝的進化 $\rightarrow$ 相互作用を伴う共進化」と進んできたが，\textbf{協調的な相互作用（議論）のチームレベル評価を適応度とする世代交代型LoRA進化}は空白のまま残されている．

\section{協調の定量化と適応度設計の理論}
\label{sec:2.4}

前節の空白を埋めるには，「チームへの貢献」を原理的に定義し，それを進化の選択圧として健全に機能させる理論が必要である．本節では本研究の適応度設計が依拠する4つの理論的支柱を整理する．

\subsection{Shapley値による貢献配分}
\label{sec:2.4.1}

協力ゲーム理論のShapley値\cite{shapley1953value}は，連合（coalition）の価値関数$v$が与えられたとき，各プレイヤーの貢献を全部分連合への限界貢献の加重平均として配分する．機械学習分野ではGhorbaniとZou\cite{ghorbani2019data}が訓練データの価値評価（Data Shapley）に導入し，Shapley値が効率性・対称性・null player・線形性の4公理を満たす一意の配分であることを配分原理としての根拠とした．Rozemberczkiら\cite{rozemberczki2022shapley}のサーベイが整理するように，Shapley値の計算は一般に$2^n$個の連合評価を要するため近似手法が発達してきたが，逆に言えば$n$が小さければ厳密計算が可能である．本研究は議論チームを3エージェントに設計することで，全7連合（空集合を除く）の議論を実測し\textbf{近似誤差のない厳密Shapley値}を得る．これは単純なleave-one-out（LOO）近似が持つ欠陥——たとえば同一能力の複製が2個体いる場合，どちらを除いても価値が変わらないため両者に貢献0を割り当ててしまう——を原理的に回避する．

マルチエージェント強化学習におけるcredit assignmentの系譜も同じ問題意識を共有する．WolpertとTumer\cite{wolpert1999collective}のdifference reward（Wonderful Life Utility）は「自分がいなかった場合」との反事実差分で個体の貢献を測る枠組みの源流であり，Foersterら\cite{foerster2018counterfactual}のCOMAはcounterfactual baselineを用いた方策勾配としてこれを深層MARLに実装した．Shapley値はこれら反事実的貢献を全連合にわたって公理的に平均化したものと位置づけられる．

\subsection{協調的共進化と集団ベース訓練}
\label{sec:2.4.2}

複数の部分要素が組み合わさって初めて全体の性能が決まる問題に対する進化計算の古典的解が，PotterとDe Jong\cite{potter1994cooperative}の協調的共進化（Cooperative Coevolution）である．問題を部分要素（本研究では議論の役割）ごとのサブ集団に分割し，各個体を他サブ集団の代表（協力者）と組み合わせて評価する．この「代表チーム文脈での個体評価」は，本研究において各候補LoRAのShapley値を計算する際のチーム構成法（候補1体＋他役割の代表2体）の骨格そのものである．また，Jaderbergら\cite{jaderberg2017population}のPopulation Based Training（PBT）は，訓練途中の集団に対するexploit（劣位個体の優位個体による置換）とexplore（摂動）の交互適用が有効であることを示しており，本研究の世代ループ（役割内選抜と突然変異）の手続き的な根拠を与える．

\subsection{多様性の扱い：加算ではなく分離・ペナルティ}
\label{sec:2.4.3}

進化計算において目的関数への単一目的化は欺瞞的（deceptive）な探索を招きうる．LehmanとStanley\cite{lehman2011abandoning}のNovelty Searchは，目的への直接最適化がかえって目的達成を妨げる場合があることを示し，行動の新規性そのものを探索圧とする発想を導入した．しかし「では多様性を適応度に加算すればよい」とは言えない．MouretとClune\cite{mouret2015illuminating}のMAP-Elitesは，多様性を適応度に混ぜるのではなく行動記述子で張られたアーカイブとして分離保持する品質多様性（QD）の枠組みを確立し，多目的最適化の古典NSGA-II\cite{deb2002fast}も，加重和方式が非凸Pareto前線を捉えられず重み係数の選択が恣意的になるという限界を明確にしている．アンサンブル学習の側からは，KroghとVedelsby\cite{krogh1995neural}のambiguity分解——アンサンブル誤差$E$は構成員の平均誤差$\bar{E}$から構成員間の不一致$\bar{A}$を引いた$E = \bar{E} - \bar{A}$に分解される——が多様性の価値の数学的根拠を与える一方，Woodら\cite{wood2023unified}の統一理論は，多様性がbias・varianceと並ぶ誤差の第3次元であり，損失関数に依存してその効き方が変わるトレードオフ管理の対象であることを示し，多様性の無条件な加算に対する最も強い反証となっている．

これらを踏まえ，本研究は多様性を適応度への加算項とせず，GoldbergとRichardsonによるfitness sharing\cite{goldberg1987genetic}——行動的に近接した個体同士で適応度を割り引く乗法ペナルティ——として組み込む．fitness sharingはニッチ形成による早期収束の防止という明確な役割を持ち，貢献度（Shapley値）の順序構造を加算項のように恣意的な係数で歪めない．またKroghとVedelsbyの分解は，進化の過程でチームの行動多様性（回答不一致プロファイル）がどう変化しチーム性能とどう関係するかを分析する枠組み（研究課題(c)の多様性設計に対応）として用いる．

\subsection{評価方法論}
\label{sec:2.4.4}

適応度と最終評価の設計には評価方法論の知見を用いる．tinyBenchmarks\cite{polo2024tinybenchmarks}は約100問のサブセットでフルベンチマークとの誤差が約2\%に収まることを示しており，進化ループ内の適応度セット（固定100問）の規模設定の根拠となる．最終評価の統計的検定は，SEMの併記・クラスタ標準誤差・複数回リサンプル・paired分析・検出力分析を勧告するMiller\cite{miller2024adding}の指針と，NLPにおける有意性検定の使い分けを整理したDrorら\cite{dror2018hitchhiker}に従う．

\section{本研究の位置づけ}
\label{sec:2.5}

以上の4系統を俯瞰すると，既存研究は「協調プロトコルの改善（\ref{sec:2.1}節）」「重み合成の演算子（\ref{sec:2.2}節）」「重み集団の進化（\ref{sec:2.3}節）」「貢献と多様性の理論（\ref{sec:2.4}節）」をそれぞれ発展させてきたが，それらの交点——\textbf{議論のチームレベル適応度を選択圧とするLoRA重み集団の世代交代型進化}——は未報告である．表\ref{tab:2.2}に，隣接する研究系統と本研究に欠けている要素を整理する．

\begin{table}[htbp]
\caption{隣接研究系統と本研究の位置づけ}
\label{tab:2.2}
\centering
\footnotesize
\setlength{\tabcolsep}{3pt}
\begin{tabular}{|p{2.4cm}|p{2.7cm}|p{2.3cm}|p{2.4cm}|p{3.4cm}|}
\hline
系統 & 代表研究 & 最適化の対象 & 最適化の信号 & 本研究との差分（欠けている要素） \\
\hline
LoRA集団＋進化演算 & GENOME\cite{zhang2025genome}, EvoPref\cite{evopref2026evolutionary} & LoRA重み集団 & 個体タスク性能 & 適応度が個体性能のみで，協調寄与を測らない \\
\hline
LoRA集団＋相互作用共進化 & PopuLoRA\cite{creuscastanyer2026populora} & LoRA重み集団 & 敵対的self-playの勝敗 & 相互作用が敵対的であり，協調議論でない \\
\hline
議論品質を最適化信号に & MAPoRL\cite{park2025maporl} & モデル重み（勾配更新） & 議論の協調品質報酬 & RLであり，集団の世代交代型進化でない \\
\hline
チーム適応度でMASを進化 & EvoMAS\cite{evomas2026evolving}, Evolve as a Team\cite{metateam2026meta} & プロンプト・エージェント構成 & チーム性能 & 進化対象が重みでなくプロンプト／構成 \\
\hline
LoRA重み空間の協調探索 & Model Swarms\cite{feng2024model} & LoRA重み集団 & 個体効用 & PSOであり交叉・突然変異・世代交代がない \\
\hline
議論由来データでSFT & Multiagent Finetuning\cite{subramaniam2025multiagent} & モデル重み（勾配更新） & 議論トレースの模倣 & 勾配による模倣であり，選択圧でない \\
\hline
\textbf{本研究} & --- & \textbf{役割別LoRA重み集団} & \textbf{議論チームへの厳密Shapley寄与 $\times$ fitness sharing} & --- \\
\hline
\end{tabular}
\end{table}

本研究の位置づけは次の3点に要約される．

第一に，\textbf{最適化信号の新規性}である．表\ref{tab:2.2}の上段2系統（GENOME・PopuLoRA）は重み集団の進化という「器」を持つが信号が個体的・敵対的であり，中段2系統（MAPoRL・EvoMAS等）は協調・チームという「信号」を持つが対象が勾配更新やプロンプトである．本研究は，公理的に正当化された貢献配分である厳密Shapley値（\ref{sec:2.4.1}項）をLoRA重み集団の世代交代進化の適応度に導入することで，この器と信号を初めて接続する．GENOME型の個体性能適応度はアブレーション（A1）として設計・実装したが，計算資源の制約から本実験では未実施である（第\ref{cha:5}章）．新規性の主張は定式化の水準に留まる点に注意されたい．

第二に，\textbf{理論的整合性を持つ演算子の選択}である．交叉はKnOTSの知見（\ref{sec:2.2}節）に基づく$\Delta W$空間ブレンド＋SVD再分解とし，因子別補間の交差項問題を回避する．多様性はQDとアンサンブル統一理論の批判（\ref{sec:2.4.3}項）を踏まえ，加算項ではなくfitness sharingの乗法ペナルティとして扱う．議論プロトコルは批判研究が特定した成功条件（\ref{sec:2.1.3}項：独立初期回答・少数ラウンド・多数決集約）に準拠する．

第三に，\textbf{否定的知見も成果となる問題設定}である．4B級の同質エージェントによる素の議論は機能しない可能性が高いという緊張関係（\ref{sec:2.1.3}項）の下で，チームレベル選択圧が議論を機能させるならばそれはMADを「エージェント集団の育成問題」として捉え直す証拠となり，機能しないならば適応度別・世代別の性能と行動多様性の軌跡から小型モデル議論の成立条件を統計的に特定できる．いずれの帰結でも，計算量を統制したSelf-Consistencyベースライン（Smitら\cite{smit2024should}の要請）とpaired検定（Miller\cite{miller2024adding}の指針）に基づく厳密な比較の上で結論が導かれる点が，評価不備が指摘されてきたMAD研究\cite{zhang2025debate}に対する方法論的な貢献となる．

なお，本領域は2025〜2026年にかけて発展が速く，PopuLoRAのような近接研究が継続的に出現している．本表は執筆時点（2026年7月）の再検索に基づく．
